\section{Limitations}

\begin{itemize}
  \item \textbf{Local students see frames, not video.} The student receives up to 4 frames per second at 512 pixels
  (about 70 tokens a frame), without audio. Code2Video's own protocol (Gemini on the mp4) is implemented and exported
  for grading by hand, but not run.
  \item \textbf{One student family.} The main student is \texttt{gemma4:26b}; the prior was also measured with
  \texttt{qwen2.5vl:7b}. Students of other families may know more or less.
  \item \textbf{Invented concepts are small.} NOVEL-12's concepts are definable on one card; they test whether a video
  transmits a definition and its consequences, not deep understanding of a large theory. TEST was written by one model
  (Claude) and not reviewed by teachers.
  \item \textbf{The evidence judge is not calibrated.} Its decisions have not been compared with human annotations
  (e.g. Cohen's $\kappa$~\cite{cohen1960}); the forced-choice versions do not use it.
  \item \textbf{Twelve topics per set.} Means over 12 topics move a lot with one topic. We do not pool questions
  across topics, because the topic, not the question, is what receives a repair; a larger set would make the means
  steadier.
  \item \textbf{The local coder.} After one retry, two of twelve NOVEL-12 videos were still incomplete and only
  \AdmClean{} of \AdmTotal{} were clean. Stronger code models are future work; the Claude-fixed versions show the
  room.
  \item \textbf{The repair is not yet evaluated on NOVEL-12.} F0--F5 are built and their rules fixed in advance; the runs continue.
  \item \textbf{The model-free checks are narrow.} The timing check covers two kinds of fault; the attachment check
  skips labels anchored to a text glyph.
  \item \textbf{Counterfactual relabelling was not applied to MMMC.} It would rewrite TEST per video and break the
  paired design; NOVEL-12 controls the prior instead.
\end{itemize}

\section{Conclusion}

TeachQuiz, as published and as we first ran it, could not have detected teaching: a grading bug dropped answers,
the student's refusals set $S_1$, and the questions were answered from prior knowledge. Forced choice, a measured
prior, normalisation, dense silent frames, a motion control, continuous scoring and option rotation each address a
named failure, and the ladder shows which of them a teaching metric actually needs. A benchmark whose prior is
controlled by construction is the precondition for asking whether a repair teaches. The layout score has the same
weakness from the other side: it measures box separation, not whether a viewer can read what belongs to what. The
repair evaluation itself (F0--F5) is the next step.
